\chapter{Planteamiento del problema}\label{cap6}

\section{Exposición del problema}

Investigadores de la ENCB se encuentran estudiando las propiedades antibióticas
de una variedad de organismos pertenecientes al género ascomicetos coprófilos,
sin embargo, el procedimiento actual en cajas Petri está limitado a una producción
reducida, además de requerir una alta manipulación manual. Existen biorreactores
comerciales para el escalado de estos sistemas biológicos, sin embargo, ninguno
considera las particularidades de estos organismos, como la baja resistencia a las
fuerzas cortantes, el alto índice de oxigenación requerido, iluminación homogénea
y control de otras variables ambientales clave (pH, presión, temperatura). Esta
brecha tecnológica representa una limitante en la reproducibilidad y eficacia de los
experimentos a efectuar sobre estos organismos.

\section{Solución propuesta}

El presente proyecto tiene como objetivo diseñar y construir un biorreactor de tipo
airlift con una capacidad útil de 43 litros, orientado a la investigación en el cultivo
controlado de hongos coprófilos del filo Ascomycota. El biorreactor será diseñado
considerando las características particulares del hongo, proporcionando un entorno
controlado y libre de contaminantes significativos. Para asegurar una repartición
consistente de los nutrientes y una oxigenación adecuada, se añadirá un sistema
de burbujeo que permita la circulación del medio. Además, se incorporarán
mecanismos para monitorear y controlar las variables temperatura, iluminación y
pH a lo largo de las diferentes etapas de crecimiento del hongo.

Se contará con un subsistema que permitirá liberar el exceso de oxígeno y otros
gases generados durante el proceso con el fin de evitar una zona de aireación en
la cual se puedan desarrollar elementos no deseados como el micelio
contaminante. El biorreactor integrará una interfaz humano-máquina que permitirá
al usuario programar el ensayo a realizar en el hongo de manera intuitiva. Durante
la prueba, mostrará los valores más recientes correspondientes al muestreo con
actualizaciones cada minuto.

Al finalizar el experimento, permitirá descargar la información capturada en una
memoria externa para su posterior análisis. Se busca que el diseño del biorreactor
sea modular, puesto que el desmontaje para limpieza y esterilización es un factor
clave en el cultivo. Todos los materiales de construcción deberán ser de grado
alimenticio y compatibles con los protocolos de limpieza química. Esta
configuración modular también facilitará la recolección de la biomasa generada,
permitiendo el acceso directo a los sólidos producidos posterior al vaciado del
contenido líquido. El sistema completo deberá tener dimensiones máximas de
80~cm $\times$ 80~cm $\times$ 100~cm, e incorporar compuertas para la adición de los reactivos.
